\scanpage{7}
% archive foliation: 6
\[\bigwedge^n M\]

a) Cas où $r=0$, $r=1$

b) Fonctorialité en $X,Y,S$
\inserted{\diagram{\begin{tikzcd}[row sep=small,column sep=small,execute at end picture={
\draw ([xshift=-8pt]\tikzcdmatrixname-1-1.north west) to[out=-90,in=90] ([xshift=-12pt]\tikzcdmatrixname-2-1.west) to[out=-90,in=150] ([xshift=-1pt]\tikzcdmatrixname-3-1.south west);
}]
\text{(i)}&X_S^{\unsure{r|n}}\arrow[r]&Y&&&\\
{}&X\arrow[r]&X'&&Y\arrow[r]&Y'\\
\text{(ii)}&F&&&F'&\\
{}&X\arrow[d]&X_S^r\arrow[r,"\varphi"]&Y&X'\arrow[d]&X_{S'}'{}^r\arrow[r,"\varphi'"]&Y'\\
{}&S&&&S'&&
\end{tikzcd}}}
\diagram{\begin{tikzcd}[row sep=small,column sep=small]
X_S^n\arrow[r,"\varphi"]\arrow[d]&Y\arrow[d]&S\arrow[l,"\unsure{\operatorname{Cor}|\operatorname{Cn}}_F^\varphi"']\\
X_S'{}^n\arrow[r,"\varphi_1"']&Y'&S\arrow[l,"\unsure{\operatorname{Cor}|\operatorname{Cn}}_{F}^{\varphi',\unsure{u}}"']
\end{tikzcd}}
\margin{\unsure{comm. polyp.}\qquad\uuline{\unsure{comm}}}
\diagram{\begin{tikzcd}[row sep=small,column sep=large]
Y&Y'\arrow[l]\\
S\arrow[u,"\unsure{\operatorname{Cor}|\operatorname{Cn}}_F^\varphi"]\arrow[r,phantom,"\text{\uuline{comm}}"]&S'\arrow[l]\arrow[u,"\unsure{\operatorname{Cor}|\operatorname{Cn}}_{F'}^{\varphi'}"']
\end{tikzcd}}

c) Additivité en $M$
\diagram{\begin{tikzcd}[row sep=large,column sep=large]
X_S^{r_1+r_2}\arrow[r,"\varphi"]&Y&S\arrow[l,"\unsure{\operatorname{Cor}|\operatorname{Cn}}_{\deleted{\unsure{F}}F_1\unsure{\oplus}F_2}^{\varphi}"']\arrow[d,equal]\\
X_S^{r_1}\times_S X_S^{r_2}\arrow[u,"\unsure{\iota}","\mathrm{can.}"']\arrow[r,"\varphi_1\times_S\varphi_2"']&Y_1\times_S Y_2\arrow[u,"u"]&\deleted{\unsure{S}}S\arrow[l,"\unsure{\operatorname{Cor}|\operatorname{Cn}}_{F_1}^{\varphi_1}\unsure{\times}\unsure{\operatorname{Cor}|\operatorname{Cn}}_{F_2}^{\varphi_2}"]
\end{tikzcd}}
\margin{\uline{comm.}\quad\deleted{\unsure{\text{polyp.}}}}

d) Transitivité.
\diagram{\begin{tikzcd}[row sep=small,column sep=small]
G&X\arrow[d]&G\otimes_Y F&\unsure{\text{Cas du}}\ \unsure{*}\\
r&Y\arrow[d]&\unsure{rs}&\\
F&Z&&\\
s&&&
\end{tikzcd}}
\diagram{\begin{tikzcd}[row sep=small,column sep=small,execute at end picture={
\draw ([xshift=-5pt,yshift=8pt]\tikzcdmatrixname-1-1.north west)--([xshift=-5pt,yshift=-5pt]\tikzcdmatrixname-4-1.south west);
\draw ([xshift=-5pt,yshift=8pt]\tikzcdmatrixname-1-1.north west)--([xshift=8pt,yshift=8pt]\tikzcdmatrixname-1-3.north east) to[out=0,in=90] ([xshift=14pt]\tikzcdmatrixname-5-3.east)--([xshift=-55pt,yshift=-8pt]\tikzcdmatrixname-5-3.south);
\draw (\tikzcdmatrixname-2-2.south) to[out=-70,in=150] (\tikzcdmatrixname-5-3.west);
\draw (\tikzcdmatrixname-1-3.south) to[out=-60,in=70] (\tikzcdmatrixname-4-3.east);
}]
X_Y^{\unsure{r}}\arrow[r,"\varphi"]&\mathcal{X}_r&\text{($Y$-\unsure{morphisme})}\\
Y_Z^s\arrow[r,"\psi"]&\mathcal{Y}_s&\text{($Z$-\unsure{morphisme})}\\
X_Z^r\arrow[r,"\unsure{f}"]&\mathcal{Z}_r&\text{($Z$-\unsure{morphisme})}\\
\mathcal{X}_Z^s\arrow[r,"u"]&\mathcal{Z}&\text{($Z$-\unsure{morphisme})}\\
&&\text{\unsure{soit comm.}}
\end{tikzcd}}
\unsure{tels que}
\diagram{\begin{tikzcd}[row sep=large,column sep=large,execute at end picture={
\draw ([xshift=-45pt,yshift=15pt]\tikzcdmatrixname-1-1.north west)--([yshift=15pt]\tikzcdmatrixname-1-1.north east);
\draw ([xshift=-45pt,yshift=15pt]\tikzcdmatrixname-1-1.north west) to[out=-85,in=70] ([xshift=-65pt,yshift=-10pt]\tikzcdmatrixname-2-1.south west) to[out=-110,in=180] ([yshift=-35pt]\tikzcdmatrixname-2-1.south) to[out=0,in=-90] ([xshift=12pt]\tikzcdmatrixname-2-2.east);
\draw ([xshift=-25pt]\tikzcdmatrixname-1-1.west) to[out=-40,in=140] (\tikzcdmatrixname-2-1.south);
}]
X_Z^{\unsure{rs}}\arrow[r,"\unsure{\gamma}"]&\mathcal{Z}\\
(X_Y^r)_Z^s\arrow[u,"\mathrm{can}"]\arrow[r,"\varphi_Z^s"']&\mathcal{X}_Z^s\arrow[u,"u"']
\end{tikzcd}}
\margin{\unsure{commutatif}}
